% 05-copy-engine.tex - Copy Engine

\chapter{Copy Engine}
\label{sec:copy-engine}

This section describes the \code{InterruptibleCopy} class that provides the file transfer capabilities of \SC{}.

\section{Overview}

The copy engine wraps \code{rsync} to provide reliable, resumable file transfers with real-time progress monitoring. Each active copy operation is managed by an \code{InterruptibleCopy} instance that runs the \code{rsync} process in a background thread and monitors its output.

\section{Copy Modes}

\SC{} supports four copy modes, determined by the relationship between the source host and destination host (Table~\ref{tab:copy-modes}).

\begin{table}[htbp]
    \centering
    \caption{Copy Modes}
    \label{tab:copy-modes}
    \begin{tabular}{lp{4cm}p{5cm}}
        \toprule
        \textbf{Mode} & \textbf{Condition} & \textbf{Mechanism} \\
        \midrule
        Local to local & source host = ``'', dest host = ``'' & \code{rsync} locally, no SSH \\
        Local to remote & source host = ``'', dest host $\neq$ ``'' & \code{rsync} from local filesystem to remote destination \\
        Remote same-host & source host = dest host $\neq$ ``'' & SSH to \DR{}, \code{rsync} locally on the \DR{} \\
        Remote to remote & source host $\neq$ ``'', source host $\neq$ dest host & SSH to source \DR{}, \code{rsync} to remote destination \\
        \bottomrule
    \end{tabular}
\end{table}

\subsection{Local to Local Copy}

For locally originating copies with a local destination, \SC{} executes:
\begin{lstlisting}
rsync -avH --append --partial --progress <src> <dest>
\end{lstlisting}

No SSH wrapper is used.

\subsection{Local to Remote Copy}

For locally originating copies with a remote destination, \SC{} executes:
\begin{lstlisting}
rsync -avH --append-verify --partial --progress <src> <dest>:<destpath>
\end{lstlisting}

Local to remote copies use \code{--append-verify} instead of \code{--append} for additional integrity checking. The bandwidth limit is applied when the destination is determined to be behind a gateway (see Section~\ref{sec:bw-gateway-detection}).

\subsection{Remote Same-Host Copy}

For copies where the source and destination are on the same \DR{}, \SC{} executes:
\begin{lstlisting}
ssh -t -t mcsdr@<DR> 'shopt -s huponexit && \
    rsync -avH --append --partial --progress <src> <dest>'
\end{lstlisting}

The \code{shopt -s huponexit} ensures the \code{rsync} process is terminated if the SSH session is disconnected, preventing orphaned processes on the \DR{}.

\subsection{Remote to Remote Copy}

For copies where the source and destination are on different hosts, \SC{} executes:
\begin{lstlisting}
ssh -t -t mcsdr@<DR> 'shopt -s huponexit && \
    rsync -avH --append-verify --partial --progress \
    [--bwlimit=<limit>m] <src> <dest>:<destpath>'
\end{lstlisting}

Remote to remote copies use \code{--append-verify} instead of \code{--append} for additional integrity checking. The bandwidth limit is applied when the destination is determined to be behind a gateway (see Section~\ref{sec:bw-gateway-detection}).

\section{Pause and Resume}

Copy operations can be paused and resumed:

\begin{description}
    \item[Pause] The \code{rsync} subprocess is killed, terminating the transfer. The copy state is preserved and the status is set to ``paused''.

    \item[Resume] A new \code{rsync} process is started with the same parameters. The \code{--append} (or \code{--append-verify}) flag ensures the transfer resumes from where it left off, only transferring bytes not already present at the destination.
\end{description}

Resuming a paused copy does not increment the retry counter.

\section{Destination File Truncation}
\label{sec:truncation}

Before resuming a copy, \SC{} truncates the last 512~KB of the destination file. This addresses a failure mode where the final block of a partially transferred file may be corrupted due to an interruption during write. Because local and same-host copies use \code{rsync --append} (which does not verify existing data), truncation ensures the corrupted tail is discarded before appending resumes.

The truncation is performed via:
\begin{lstlisting}
truncate -c -s -512K <destination_file>
\end{lstlisting}

This operation is applied only for local and same-host copies. Remote copies use \code{rsync --append-verify}, which checksums existing data, so truncation is not needed.

\section{Bandwidth Limiting}
\label{sec:bw-gateway-detection}

For remote copies where the \code{bw\_limit} configuration parameter is non-zero, \SC{} applies \code{rsync}'s \code{--bwlimit} option to limit the transfer rate. Bandwidth limiting is only applied when the destination host is determined to be behind a network gateway, indicating that the traffic must traverse a wide-area link. This gateway-based definition of ``remote'' is stricter than the hostname comparison used for copy exclusivity (Section~\ref{sec:remote-exclusivity}): a copy to a different host on the local network is exclusive but not bandwidth-limited.

\subsection{Gateway Detection}

Before starting a copy, \SC{} resolves the destination hostname to an IP address and queries the kernel routing table using \code{ip route get}. If the route to the destination includes a gateway hop (indicated by \code{via} in the output), the destination is considered remote and bandwidth limiting is applied.

For locally originating copies, the route check runs on the local host:
\begin{lstlisting}
ip route get <dest_ip>
\end{lstlisting}

For \DR{}-originating copies, the route check runs on the \DR{} via SSH:
\begin{lstlisting}
ssh -t -t mcsdr@<DR> 'ip route get <dest_ip>'
\end{lstlisting}

If DNS resolution fails or the route check cannot be performed, the destination is assumed to be remote and bandwidth limiting is applied.

\section{Progress Monitoring}

The copy engine monitors \code{rsync} output in real time by parsing the progress line format:

\begin{lstlisting}
<bytes_transferred> <percent>% <speed> <time_remaining>
\end{lstlisting}

The following progress metrics are extracted:
\begin{itemize}
    \item \textbf{Bytes Transferred}: Total bytes written to the destination
    \item \textbf{Progress}: Percentage complete (e.g., ``42\%'')
    \item \textbf{Speed}: Current transfer rate (e.g., ``12.34MB/s'')
    \item \textbf{Time Remaining}: Estimated time to completion (e.g., ``01:23:45'')
\end{itemize}

These values are available via the per-\DR{} \mibentry{ACTIVE\_\{property\}\_DR\{n\}} MIB entries (Section~\ref{sec:mib-entries}).

\section{Delete Operations}

Delete operations use \code{rm -f} rather than \code{rsync}:

\begin{lstlisting}
ssh -t -t mcsdr@<DR> 'shopt -s huponexit && sudo rm -f <path>'
\end{lstlisting}

For locally originating deletes:
\begin{lstlisting}
rm -f <path>
\end{lstlisting}

\section{Copy Lifecycle}

The lifecycle of a copy operation is:

\begin{enumerate}
    \item \textbf{Creation}: An \code{InterruptibleCopy} instance is created with the source, destination, ID, retry count, and last try timestamp
    \item \textbf{Retry Check}: If the last attempt was less than \code{wait\_retry} hours ago (default: 24), the copy sets its status to ``error: too soon to retry'' and does not start
    \item \textbf{Start}: The \code{rsync} process is launched in a background thread
    \item \textbf{Monitoring}: The main thread polls \code{rsync} stdout/stderr for progress updates
    \item \textbf{Completion}: When \code{rsync} exits:
    \begin{itemize}
        \item Return code 0: status set to ``complete''
        \item Return code $<$ 0 (killed by signal): status set to ``paused''
        \item Return code $>$ 0: status set to ``error: \textit{stderr output}'', and any rsync error messages from stdout are prepended
    \end{itemize}
    \item \textbf{Post-Processing}: The queue manager handles the result (log completion, schedule retry, or record failure)
\end{enumerate}
